\documentclass[12pt,a4paper]{article}

\usepackage[utf8]{inputenc}
\usepackage[T1]{fontenc}
\usepackage[brazil]{babel}
\usepackage{amsmath,amssymb,amsthm}
\usepackage{proof}          % \infer
\usepackage{bussproofs}     % árvores de prova estilo Gentzen
\usepackage{enumitem}
\usepackage{geometry}
\usepackage{listings}
\usepackage{xcolor}
\usepackage{hyperref}
\usepackage{mathtools}

\geometry{margin=2.5cm}

% ── código Haskell ────────────────────────────────────────────────────────────
\lstdefinelanguage{Haskell}{
  keywords={data,type,where,let,in,case,of,do,module,import,if,then,else,
  deriving,newtype,class,instance,Right,Left,Just,Nothing},
  keywordstyle=\bfseries\color{blue!70!black},
  comment=[l]{--},
  commentstyle=\itshape\color{gray},
  stringstyle=\color{orange!80!black},
  basicstyle=\ttfamily\small,
  breaklines=true,
  showstringspaces=false
}
\lstset{language=Haskell}

% ── macros lógicas ────────────────────────────────────────────────────────────
\newcommand{\seq}{\;\vdash\;}
\newcommand{\seqc}[2]{#1\vdash #2}
\newcommand{\arr}{\supset}
\newcommand{\lc}{\lambda}
\newcommand{\FV}{\mathsf{FV}}
\newcommand{\Red}{\mathbf{Red}}

% ── contadores ────────────────────────────────────────────────────────────────
\newcounter{exercicio}
\newenvironment{exercicio}[1][]{%
  \refstepcounter{exercicio}%
  \medskip\noindent\textbf{Exercício \theexercicio%
  \ifx\relax#1\relax\else\ (#1)\fi.}\enspace\ignorespaces
}{%
  \par\smallskip
}

% ─────────────────────────────────────────────────────────────────────────────
\title{\textbf{Lógica Aplicada à Computação}\\[4pt]
\large Lista de Exercícios 1\\[2pt]}
\author{Prof.\ Rodrigo Ribeiro}
\date{}

\begin{document}
\maketitle
\thispagestyle{empty}

\bigskip
\noindent\textit{Instruções gerais:}
\begin{itemize}[leftmargin=1.5em]
  \item Apresente suas respostas de forma clara, identificando as regras de
    inferência utilizadas.
  \item Em exercícios de redução, mostre cada passo.
  \item Entrega deverá ser feita exclusivamente pelo Moodle.
  \item Sua solução deverá ser submetida como um único arquivo PDF,
    produzido utilizando~\LaTeX.
\end{itemize}

% ─────────────────────────────────────────────────────────────────────────────
\section*{Parte I — Lógica Proposicional: Dedução Natural}
% ─────────────────────────────────────────────────────────────────────────────

\begin{exercicio}
  Construa derivações em dedução natural para cada um dos sequentes abaixo.
  Identifique explicitamente cada regra aplicada.
  \begin{enumerate}[label=(\alph*)]
    \item $\seqc{}{A \arr (B \arr A)}$
    \item $\seqc{A \arr B,\; B \arr C}{A \arr C}$
    \item $\seqc{}{(A \land B) \arr (B \land A)}$
    \item $\seqc{}{A \arr (A \lor B)}$
    \item $\seqc{}{(A \lor B) \arr (\neg A \arr B)}$
    \item $\seqc{\neg A \land \neg B}{\neg(A \lor B)}$
  \end{enumerate}
\end{exercicio}

\section*{Parte II — Lógica Proposicional: Cálculo de Sequentes}

\begin{exercicio}
  Utilize o sistema de proof search (Dyckhoff, 1992) para construir
  provas dos seguintes
  sequentes ou demonstre que eles \emph{não} são demonstráveis na lógica
  intuicionista.  Em cada caso, exiba a árvore de prova completa ou a indicação
  de falha.
  \begin{enumerate}[label=(\alph*)]
    \item $\seq P \arr P$
    \item $\seq (P \land Q) \arr (Q \land P)$
    \item $\seq (P \arr Q \arr R) \arr (P \arr Q) \arr P \arr R$
      \hfill\textit{(combinador S)}
    \item $P \arr Q,\; P \seq Q$
    \item $\seq P \lor \neg P$ \hfill\textit{(terceiro excluído)}
    \item $\seq \neg\neg P \arr P$ \hfill\textit{(dupla negação)}
  \end{enumerate}
\end{exercicio}

\section*{Parte III — Lambda Cálculo Não Tipado}

\begin{exercicio}[Notação De Bruijn]
  Converta os seguintes termos nomeados para a notação De Bruijn.
  \begin{enumerate}[label=(\alph*)]
    \item $\lc x.\; \lc y.\; x$
    \item $\lc x.\; \lc y.\; y$
    \item $\lc f.\; \lc x.\; f\;(f\;(f\;x))$
    \item $(\lc x.\; x\; x)\; (\lc x.\; x\; x)$
      \hfill\textit{(combinador $\Omega$)}
    \item $\lc x.\; \lc y.\; \lc z.\; x\;z\;(y\;z)$
      \hfill\textit{(combinador S)}
  \end{enumerate}
\end{exercicio}

\begin{exercicio}[Beta-redução]
  Reduza os termos abaixo à forma normal (se existir), mostrando cada passo.
  Indique qual estratégia de avaliação foi utilizada (\emph{call-by-value} ou
  \emph{normal-order}).
  \begin{enumerate}[label=(\alph*)]
    \item $(\lc x.\; x)\; (\lc y.\; y)$
    \item $(\lc x.\; \lc y.\; x)\; M\; N$ para termos quaisquer $M$, $N$.
    \item $(\lc x.\; x\; x)\; (\lc y.\; y)$
    \item $\mathbf{succ}\;\mathbf{2}$, onde
      $\mathbf{succ} = \lc n.\lc f.\lc x.\; f\;(n\;f\;x)$ e
      $\mathbf{2} = \lc f.\lc x.\; f\;(f\;x)$.
    \item $\mathbf{plus}\;\mathbf{1}\;\mathbf{2}$, onde
      $\mathbf{plus} = \lc m.\lc n.\lc f.\lc x.\; m\;f\;(n\;f\;x)$.
  \end{enumerate}
\end{exercicio}

\begin{exercicio}[Codificações de Church]
  Utilizando as codificações de Church para booleanos e números naturais,
  defina os seguintes termos e verifique a corretude através de reduções.
  \begin{enumerate}[label=(\alph*)]
    \item Defina $\mathbf{xor}$ e mostre que
      $\mathbf{xor}\;\mathbf{true}\;\mathbf{false} \to^* \mathbf{true}$.
    \item Defina $\mathbf{mult}$ (multiplicação) e mostre que
      $\mathbf{mult}\;\mathbf{2}\;\mathbf{3} \to^* \mathbf{6}$.
    \item Defina um operador de par $\mathbf{pair}$, $\mathbf{fst}$ e
      $\mathbf{snd}$ e verifique que
      $\mathbf{fst}\;(\mathbf{pair}\;M\;N) \to^* M$.
  \end{enumerate}
\end{exercicio}

\begin{exercicio}[Shift e substituição em De Bruijn]
  Seja o operador $\uparrow^d_c$ de \emph{shift} e a substituição
  $[j \mapsto s]$ na notação De Bruijn.
  \begin{enumerate}[label=(\alph*)]
    \item Calcule $\uparrow^1_0\;(\lc.\;\lc.\;1\;0)$.
    \item Calcule $[0 \mapsto \lc.\;0]\;(\lc.\;1\;0)$.
    \item Realize a beta-redução $(\lc.\;\lc.\;1\;0)\;(\lc.\;0)$ usando a
      fórmula
      $(\lc.\;t_1)\;t_2 \longrightarrow
      \uparrow^{-1}_0([0\mapsto\uparrow^1_0 t_2]\,t_1)$
      e verifique o resultado convertendo para a notação nomeada.
  \end{enumerate}
\end{exercicio}

% ─────────────────────────────────────────────────────────────────────────────
\section*{Parte IV — Lambda Cálculo Tipado Simples e Curry-Howard}
% ─────────────────────────────────────────────────────────────────────────────

\begin{exercicio}[Derivações de tipo]
  Para cada termo abaixo, construa a derivação de tipo no STLC ou explique por
  que o termo não é tipável.
  \begin{enumerate}[label=(\alph*)]
    \item $\lc x{:}A.\; x$
    \item $\lc x{:}A.\; \lc y{:}B.\; x$
    \item $\lc f{:}A\to B.\; \lc g{:}B\to C.\; \lc x{:}A.\; g\;(f\;x)$
    \item $\lc h{:}A \times B.\; (\mathbf{snd}\;h,\;\mathbf{fst}\;h)$
    \item $\lc p{:}A \times B.\; \lc q{:}C.\;
      \mathbf{inl}\;(\mathbf{fst}\;p) \;\mathbf{as}\; A + C$
    \item $\lc x{:}\bot.\; \mathbf{absurd}\;x \;\mathbf{as}\; A$
  \end{enumerate}
\end{exercicio}

\begin{exercicio}[Isomorfismo de Curry-Howard]
  Para cada proposição da lógica proposicional intuicionista, escreva o tipo
  correspondente no STLC e exiba um habitante (isto é, um termo fechado
  bem-tipado nesse tipo).  Se a proposição \emph{não} for intuicionisticamente
  válida, argumente por que o tipo não tem habitante.
  \begin{enumerate}[label=(\alph*)]
    \item $A \arr A$
    \item $(A \arr B) \arr (B \arr C) \arr A \arr C$
    \item $A \arr A \lor B$
    \item $A \land B \arr A$
    \item $(A \arr B \arr C) \arr (A \arr B) \arr A \arr C$
      \hfill\textit{(combinador S tipado)}
    \item $A \lor \neg A$ \hfill\textit{(terceiro excluído)}
    \item $\neg\neg A \arr A$ \hfill\textit{(eliminação da dupla negação)}
  \end{enumerate}
\end{exercicio}

\section*{Parte V — Lógica de Predicados}

\begin{exercicio}[Sintaxe e variáveis]
  Seja $\Sigma = \{0 : 0,\; s : 1,\; + : 2\}$ e
  $\mathcal{P} = \{\leq : 2,\; \text{Par} : 1\}$.
  \begin{enumerate}[label=(\alph*)]
    \item Quais dos seguintes são termos bem formados sobre $\Sigma$?
      \begin{enumerate}[label=(\roman*)]
        \item $s(0)$
        \item $+(x, s(0))$
        \item $s(s(+))$
      \end{enumerate}
    \item Calcule $\FV(\varphi)$ para as fórmulas:
      \begin{enumerate}[label=(\roman*)]
        \item $\forall x.\; \exists y.\; x + y = 0$
        \item $(\exists y.\; x + x = y) \lor (\forall x.\; x = z)$
      \end{enumerate}
    \item Calcule $[x \mapsto s(z)]\,(\exists y.\; x + y = 0)$.  Há captura
      de variável?  Justifique.
  \end{enumerate}
\end{exercicio}

\begin{exercicio}[Dedução natural com quantificadores]
  Construa derivações em dedução natural para os sequentes abaixo.
  \begin{enumerate}[label=(\alph*)]
    \item $\seqc{}{\forall x.\; P(x) \arr P(x)}$
    \item $\seqc{\forall x.\; P(x) \arr Q(x),\;\forall x.\;
      P(x)}{\forall x.\; Q(x)}$
    \item $\seqc{\forall x.\; P(x) \land Q(x)}{\forall x.\; P(x)}$
    \item $\seqc{\exists x.\; P(x) \land Q(x)}{\exists x.\; P(x)}$
    \item $\seqc{}{\forall x.\; \neg P(x) \arr \neg\exists x.\; P(x)}$
      \hfill\textit{(do slide da aula 06)}
  \end{enumerate}
\end{exercicio}

\begin{exercicio}[Semântica]
  Considere a estrutura $\mathcal{M} = (\mathbb{N}, \{+, \cdot, 0, 1\},
  \{\leq\})$.
  \begin{enumerate}[label=(\alph*)]
    \item Determine se cada fórmula abaixo é satisfeita em $\mathcal{M}$:
      \begin{enumerate}[label=(\roman*)]
        \item $\forall x.\; \exists y.\; x + y = 0$
        \item $\exists x.\; \forall y.\; x \leq y$
        \item $\forall x.\; \forall y.\; x \leq y \lor y \leq x$
      \end{enumerate}
    \item Para cada fórmula falsa em (a), exiba uma estrutura em que ela seja
      verdadeira.
  \end{enumerate}
\end{exercicio}
\end{document}
